\section{结论：阅读编译系统的三条原则}
\label{sec:conclusion}

本文从一个只有几十行的 SysY 截断点积出发，追踪同一项计算如何不断改变形状：宏常量消失在记号流中，表达式获得名字绑定和类型，循环成为控制流图中的回边与静态单赋值形式（Static Single Assignment, SSA）中的递归值，数组下标成为地址运算，函数调用成为寄存器约定、符号和重定位，结构化归约又可以在多层中间表示（Multi-Level Intermediate Representation, MLIR）中延后展开。贯穿这些变化的并不是某种共同的文本形状，而是各层之间对程序意义的承接关系。由此可以提炼出三条彼此依赖的原则。

\begin{enumerate}
  \item \textbf{意义是关系，表示是选择。}
  相邻表示不必相像：源语言中的循环、LLVM IR 中的基本块、RV64 中的分支序列可以实现同一个状态变换。语义关系划定哪些转换是合法的，表示选择决定哪些事实容易被分析，而成本模型才在合法候选中决定采用哪一个。把三者混为一谈，就会把“能够证明正确”误写成“值得生成”，或把“机器上更快”误写成“语义上更正确”。本案例中的逐项截断不能移到求和之后，这是意义施加的硬约束；是否展开循环、使用分支还是条件选择、是否保留压缩指令空间，则取决于代码尺寸、吞吐量、编译时间和目标微架构。正确性决定可行域，工程成本决定落点。

  \item \textbf{表示管理着信息的寿命，也管理着承诺的时机。}
  每一种中间表示都像一份信息预算：带类型抽象语法树保留源级结构和诊断位置，SSA 暴露定义--使用关系，机器表示显式化寄存器类别和指令约束，ELF 则保留尚未求值的符号关系。降低不是把同一内容换一种拼写，而是逐步提交决定，并常常不可逆地丢弃结构。过早把归约展开成地址与跳转，会让后续阶段重新猜测并行性；无限期保留高层抽象，又会扩大语义接口、合法化规则与维护成本。因此更稳健的原则不是“层次越高越好”，而是让一项信息活到最后一个能够利用它的阶段，随后明确地将它变为较低层的约束。MLIR 的方言共存把这种信息寿命从偶然的流水线顺序提升为可设计的系统结构。

  \item \textbf{局部最优必须服从跨边界协作。}
  单个函数若独占整个世界，可以自由选择寄存器、布局和调用序列；真实系统却要求不同语言、编译单元、库与装载环境共同工作。应用二进制接口（Application Binary Interface, ABI）牺牲一部分局部自由，以换取独立编译；符号和重定位推迟地址承诺，以换取归档复用、布局选择和装载灵活性；链接时优化和即时编译则用更多全局信息换取更长的构建时间、更复杂的部署边界或运行期开销。这里不存在无代价的“全局最优”：可见范围越大，潜在优化越多，但缓存、增量构建、模块自治和生态兼容性也越难维持。编译系统的成熟之处，正在于知道哪些决定值得跨越边界协调，哪些决定必须由稳定接口隔离。
\end{enumerate}

这三条原则把形式理论与工程现实连接起来。形式语义、控制流图、不动点分析、活跃性和重定位方程提供了可以推理的骨架；启发式搜索、分阶段编译、稳定接口和渐进式降低则回答如何在有限时间、有限内存和长期维护中实现它。理论不是现实编译器的简化替身，工程也不是对理论的不严谨妥协：前者说明不可越过的边界，后者决定在边界之内把复杂性放在哪里。

因此，阅读一个陌生编译系统时，不妨从一个具体对象开始：指出一个记号、AST 节点、SSA 值、机器指令、符号或重定位，问它表示哪项程序意义，在什么关系下承接上一层，显式保存了哪些事实，又已经提交了哪些决定；继而追问，这个决定是由正确性强制的，还是由成本模型偏好的，它又为独立编译、目标适配或后续优化留下了多少自由。能够回答这些问题，便不再只是沿着流水线背诵阶段名称，而是在阅读一套关于意义、信息与选择的计算系统。
